\begin{center}
{\large \textbf{АНОТАЦІЯ}}
\end{center}

\textit{Мусієнко О. О.} Формування інформаційно-комунікативної компетентності в процесі фахової підготовки майбутніх дизайнерів.~-- Кваліфікаційна наукова праця на правах рукопису.

Дисертація на здобуття наукового ступеня доктора філософії за спеціальністю 015~-- Професійна освіта (Цифрові технології).~-- Криворізький державний педагогічний університет, Кривий Ріг, 2026.\\

Дисертацію присвячено теоретичному обґрунтуванню та розробці моделі формування інформаційно-комунікативної компетентності (ІКК) в процесі фахової підготовки майбутніх дизайнерів. Актуальність дослідження зумовлена стрімкою ШІ-трансформацією дизайнерської практики (6,4-кратне зростання публікацій з генеративного ШІ в дизайн-освіті після 2022~р.), парадигмальним зсувом від <<дизайнера-виробника>> до <<дизайнера-куратора>> та системним відставанням програм професійної підготовки від вимог ринку праці.

Мета дослідження~-- теоретично обґрунтувати та розробити модель формування інформаційно-комунікативної компетентності в процесі фахової підготовки майбутніх дизайнерів на основі даних, включаючи конкретні рекомендації щодо трансформації навчальних програм.

Методи дослідження: аналіз і систематизація наукових джерел, оглядове дослідження (scoping review, 156 наукових праць, Web of Science), комплексний обчислювальний аналіз даних НАЗЯВО (122 програми, 3\,043 освітні компоненти), NLP-методи (тематичне моделювання BERTopic, семантична подібність, мережевий аналіз~--- 1\,240 вузлів, 7\,524 ребер, 10 спільнот Лувена), аналіз ШІ-розривів, анкетування (32~випускники, 9~роботодавців, 61~студент), статистичні методи (критерій Краскела--Уолліса), порівняльний аналіз з міжнародним бенчмаркінгом (7 програм провідних університетів), моделювання.

На основі оглядового дослідження визначено ІКК як інтегративну характеристику особистості дизайнера, що відображає його здатність ефективно здійснювати пошук, аналіз та використання професійної інформації (у~т.\,ч. згенерованої ШІ), комунікувати за допомогою цифрових засобів, застосовувати ШІ-інструменти дизайну та рефлексувати щодо етичних аспектів використання технологій. Обґрунтовано 4-компонентну структуру ІКК: інформаційно-аналітичний, комунікативний, технологічний та рефлексивний компоненти. Розроблено інтегративну теоретичну рамку на основі TPACK (адаптованого для дизайну), студійної педагогіки та компетентнісного підходу. Обчислювальний аналіз 122 українських програм спеціальності 022~Дизайн виявив системну прогалину: 0\% покриття компонентів <<Основи ШІ>>, <<Генеративний ШІ>> та <<Комп'ютерний зір>> при загальному рівні ШІ-інтеграції 35,7\%. Кейс-стаді трьох програм КДПУ з тріангуляцією 6~джерел даних підтвердив нисхідний тренд ІКК-змісту ($-$35\% за 2018--2021~рр.) та незадоволеність випускників ІКТ-підготовкою (до 33,3\%). Розроблено структурно-функціональну модель формування ІКК з 5~блоків (мотиваційно-цільовий, теоретико-методологічний, змістовий, організаційно-технологічний, діагностично-результативний), інтегровану модель 240-кредитного навчального плану з наскрізною ШІ-інтеграцією (4 нові ШІ-дисципліни, 76,7\% дисциплін із визначеним рівнем ШІ-інтеграції) та обґрунтовано 4 педагогічні умови формування ІКК.

Наукова новизна: \textit{уперше} операціоналізовано ІКК для дизайнерів у контексті ШІ з 4-компонентною структурою; проведено комплексний обчислювальний аналіз усіх 122 українських програм дизайну на предмет покриття ІКК із застосуванням тематичного моделювання, мережевого аналізу та аналізу розривів; розроблено модель трансформації навчальних програм на основі даних, що визначає інтеграцію ШІ для кожного типу навчальних дисциплін. \textit{Удосконалено} підходи до аналізу освітніх програм дизайну за допомогою обчислювальних методів у загальнонаціональному масштабі; діагностичні критерії сформованості ІКК дизайнерів з урахуванням ШІ-виміру; розуміння розвитку компетентностей у контексті українсько-європейської інтеграції. \textit{Набуло подальшого розвитку} компетентнісний підхід у дизайн-освіті щодо інтеграції ШІ-інструментів; адаптація моделі TPACK для контексту <<дизайн--ШІ>>; підходи до модернізації навчальних програм на основі аналізу великого масштабу.

Практичне значення: модель трансформації навчальних програм з рекомендаціями для 6~типів дисциплін; інтегрований 240-кредитний навчальний план із 4~новими ШІ-дисциплінами (<<Цифрова грамотність та основи AI>>, <<Генеративний дизайн та AI>>, <<Взаємодія людина-AI>>, <<Критичний дизайн та етика AI>>); інструменти анкетування для діагностики ІКК; аналітичний програмний інструментарій (Python/Playwright) для аналізу освітніх програм інших спеціальностей.

\textbf{Ключові слова}: інформаційно-комунікативна компетентність, дизайн-освіта, штучний інтелект, генеративний ШІ, професійна підготовка, навчальні програми, TPACK, компетентнісний підхід, трансформація курикулуму, prompt engineering.

\begin{center}
{\large \textbf{SUMMARY}}
\end{center}

\textit{Musiienko O. O.} Formation of Information-Communicative Competence in the Professional Training of Future Designers.~-- Qualifying scientific work as a manuscript.

Dissertation for the degree of Doctor of Philosophy (PhD) in specialty 015~-- Professional Education (Digital Technologies).~-- Kryvyi Rih State Pedagogical University, Kryvyi Rih, 2026.\\

The dissertation is devoted to the theoretical substantiation and development of a model for the formation of information-communicative competence (ICC) in the professional training of future designers. The relevance of the research is determined by the rapid AI transformation of design practice (6.4-fold increase in publications on generative AI in design education after 2022), the paradigmatic shift from ``designer-as-maker'' to ``designer-as-curator,'' and the systemic lag of professional training programs behind labour market requirements.

The aim of the study is to theoretically substantiate and develop a data-driven model for the formation of information-communicative competence in the professional training of future designers, including specific recommendations for curriculum transformation.

Research methods: analysis and systematization of scientific sources, scoping review (156 studies, Web of Science), comprehensive computational analysis of NAQA data (122 programs, 3,043 educational components), NLP methods (BERTopic topic modeling, semantic similarity, network analysis~--- 1,240 nodes, 7,524 edges, 10 Louvain communities), AI gap analysis, surveys (32~graduates, 9~employers, 61~students), statistical methods (Kruskal--Wallis test), comparative analysis with international benchmarking (7 programs of leading universities), modeling.

Based on the scoping review, ICC was defined as an integrative personal characteristic of a designer reflecting the ability to effectively search, analyze, and use professional information (including AI-generated content), communicate via digital means, apply AI design tools, and reflect on the ethical aspects of technology use. A 4-component ICC structure was substantiated: information-analytical, communicative, technological, and reflective components. An integrative theoretical framework was developed based on TPACK (adapted for design), studio pedagogy, and the competency approach. Computational analysis of all 122 Ukrainian programs in specialty 022~Design revealed a systemic gap: 0\% coverage of ``AI fundamentals,'' ``Generative AI,'' and ``Computer vision'' components, with overall AI integration at only 35.7\%. A case study of three KRSPU programs with triangulation of 6~data sources confirmed a declining ICC content trend ($-$35\% from 2018 to 2021) and graduate dissatisfaction with ICT training (up to 33.3\%). A structural-functional model for ICC formation comprising 5~blocks (motivational-goal, theoretical-methodological, content, organizational-technological, diagnostic-resultative) was developed, along with an integrated 240-ECTS curriculum model with cross-cutting AI integration (4 new AI courses, 76.7\% of courses with a defined AI integration level), and 4~pedagogical conditions for ICC formation were substantiated.

Scientific novelty: \textit{for the first time}, ICC was operationalized for designers in the AI context with a 4-component structure; a comprehensive computational analysis of all 122 Ukrainian design programs for ICC coverage was conducted using topic modeling, network analysis, and gap analysis; a data-driven curriculum transformation model specifying AI integration for each course type was developed. \textit{Improved}: approaches to analyzing design educational programs using computational methods at the national level; diagnostic criteria for assessing ICC formation in designers with regard to the AI dimension; understanding of competency development in the context of Ukrainian-European integration. \textit{Further developed}: the competency approach in design education regarding AI tool integration; adaptation of the TPACK model for the design-AI context; approaches to curriculum modernization based on large-scale program analysis.

Practical significance: a curriculum transformation model with recommendations for 6~course types; an integrated 240-ECTS curriculum with 4~new AI courses (``Digital Literacy and AI Fundamentals,'' ``Generative Design and AI,'' ``Human-AI Interaction,'' ``Critical Design and AI Ethics''); survey instruments for ICC diagnostics; analytical software toolkit (Python/Playwright) applicable to the analysis of educational programs in other specialties.

\textbf{Keywords}: information-communicative competence, design education, artificial intelligence, generative AI, professional training, curriculum, TPACK, competency approach, curriculum transformation, prompt engineering.

% Список публікацій здобувача за темою дисертації (§4 Наказу МОН №40)
\printpratsi{}
